\documentclass[11pt,fleqn]{article}
%%%%%%%%%%%%%%%%%%%%%%%%%%%%%%%%%%%%%%%%%%%%%%%%%%%%%%%%%%%%%%%%%%%%%%%%%%%%%%%%%%%%%%%%%%%%%%%%%%%%%%%%%%%%%%%%%%%%%%%%%%%%%%%%%%%%%%%%%%%%%%%%%%%%%%%%%%%%%%%%%%%%%%%%%%%%%%%%%%%%%%%%%%%%%%%%%%%%%%%%%%%%%%%%%%%%%%%%%%%%%%%%%%%%%%%%%%%%%%%%%%%%%%%%%%%%
\usepackage{graphicx}
\usepackage{amsfonts}
\usepackage{amssymb}
\usepackage{indentfirst}
\usepackage{graphics}
\usepackage{amsmath}
\usepackage{tikz}
\usepackage{graphicx}
\usepackage{hyperref}
\tikzstyle{ov}=[shape=rectangle,
                draw=black!50,
                thick,
                minimum width=0.7cm,
                minimum height=0.7cm]

\tikzstyle{av}=[shape=rectangle,
                draw=black!50,
                fill=black!10,
                thick,
                minimum width=0.7cm,
                minimum height=0.7cm]

\tikzstyle{lv}=[shape=circle,draw=black!50,thick]

\usetikzlibrary{shapes,calendar,matrix,backgrounds,folding,snakes}
\setlength{\topmargin}{-1cm}
\setlength{\oddsidemargin}{0cm}
\setlength{\evensidemargin}{0cm}
\setlength{\textheight}{22.5cm}
\setlength{\textwidth}{15.5cm}
\setlength{\parindent}{.5cm}
\newcommand{\sn}{signifcantieniveau}
\newcommand{\e}{\mbox{E}}
\newcommand{\cum}{\mbox{cum}}
\newcommand{\plim}{\mbox{plim}}
\newcommand{\var}{\mbox{var}}
\newcommand{\hvar}{\hat{\mbox{var}}}
\newcommand{\cov}{\mbox{cov}}
\newcommand{\ir}{\mbox{ir}}
\newcommand{\expit}{\mbox{expit}}
\newcommand{\logit}{\mbox{logit}}
\newcommand{\hir}{\hat{\mbox{{\i}r}}}
\newcommand{\ur}{\mbox{ur}}
\newcommand{\pr}{\mbox{pr}}
\newcommand{\cip}{\mbox{$\perp\!\!\!\perp$}}
\newcommand{\indep}{\rotatebox[origin=c]{90}{$\models$}}

\begin{document}

\title{\vspace*{-2.5cm} {\LARGE {\bf Project: Causal Machine Learning}} \\[1ex]
Stijn Vansteelandt and Federico Bertoia\\[1ex]
Academic year 2025-2026}
\date{}
\maketitle

\vspace*{-1.5cm}\\
\hrule
\vspace*{0.5cm}

\section{Introduction}

We consider estimation of the effect of eligibility and participation in 401(k) (this is a tax-advantaged retirement savings plan) on accumulated assets using 9\,915 observations at the household level, drawn from the Survey of Income and Program Participation:
\begin{verbatim}
library(hdm)
data(pension)
\end{verbatim}
In this project, each data analysis question must be based on debiased/targeted machine learning (unless otherwise stated). Make sure to use 5- or 10-fold cross-fitting in all such analyses, and to base these either on a library of complementary algorithms, or a single, but carefully tuned algorithm. In all data analyses, you are expected to report an estimate and 95\% confidence interval, if possible, and to interpret the results in an accessible manner. In addition, you are expected to make some checks to evaluate and report the quality of the results (e.g., you may want to check the stability of inverse probability weights, or whether some individuals have extreme influence functions), even when this is not explicitly asked.

\section{Task}

The key problem in determining the effect of participation in 401(k) plans on accumulated assets is the fact that the decision to enroll in a 401(k) is non-random. It is generally recognized that some people have a higher preference for saving than others. It also seems likely that those individuals with high unobserved preference for saving would be most likely to choose to participate in tax-advantaged retirement savings plans and would tend to have otherwise high amounts of accumulated assets. The presence of unobserved savings preferences with these properties then implies that conventional estimates that do not account for confounding will be biased upward, tending to overstate the savings effects of 401(k) participation. 

We will first focus on {\bf the effect of eligibility for enrolling in a 401(k) plan}. Around the time 401(k)’s initially became available, people were unlikely to be basing their employment decisions on whether an employer offered a 401(k) but would instead focus on income. Thus, eligibility for a 401(k) could be taken as exchangeable conditional on income, and the causal effect of 401(k) eligibility could be directly estimated by appropriate comparison across eligible and ineligible individuals with the same income. 
\begin{enumerate}
\item Use causal forests to estimate the effect of eligibility for enrolling in a 401(k) plan on net financial assets, in function of age, income, family size, years of education, a married indicator, a two-earner status indicator, a defined benefit pension status indicator, an IRA (individual retirement account) participation indicator, and a home ownership indicator. 
\begin{enumerate}
\item Provide insight into the results and their quality. 
\item It was previously stated that eligibility for a 401(k) could be taken as exchangeable conditional on income. Is it problem that the analysis additionally conditions on other variables, e.g. age, ...? Why (not)? You may consider the use of a causal diagram to support your reasoning.
\end{enumerate}

\item Use \texttt{ntile} to create 4 subgroups based on the obtained predictions from the causal forest and estimate the average effect of eligibility for enrolling in a 401(k) plan on net financial assets in each subgroup, along with 95\% confidence intervals. % NOTE: They need to use AIPW or TMLE here.
\item Assess if there is evidence of effect heterogeneity by estimating 
\[\sigma^2=\mbox{\rm Var}\left\{E(Y^1-Y^0|L)\right\},\]
with $L$ the variables used in the causal forest. For this, use that the efficient influence function of $\sigma^2$ is 
\begin{eqnarray*}
 &&\left\{E(Y^1-Y^0|L)-E(Y^1-Y^0)\right\}^2-\sigma^2\\
 &&+2\left\{E(Y^1-Y^0|L)-E(Y^1-Y^0)\right\}\left\{\frac{A}{P(A=1|L)}-\frac{1-A}{P(A=0|L)}\right\}\left\{Y-E(Y|A,L)\right\}
\end{eqnarray*}
Explain how you performed the analysis, report an estimate of $\sigma^2$ and corresponding 95\% confidence interval, and interpret the results.
\item Visualize the average effect of eligibility for enrolling in a 401(k) plan on net financial assets in function of age, as obtained by using an orthogonal learner, but still controlling for the same variables as before. Note that the function of causal forests may not directly allow you to do that, and that you may need to use a different orthogonal learner discussed in class. 
\begin{enumerate}
\item Report a clear visual of the results.
\item Clarify if the chosen learner makes specific assumptions that were not made by the earlier causal forest. 
\end{enumerate}

\end{enumerate}

Make a \textbf{concise, but complete} report describing what analyses you performed
and what are your conclusions to the various questions. There is no need to be extensive in your report. For instance, there is no need to repeat the theory that was used, as you may assume that the reader is familiar with the procedures that you used. However, you should carefully interpret the reported estimates and make
sure that their meaning is clear and not obscure. 
%Where possible, try to explain the differences observed between estimates obtained via different methods and give insight to the reader. {\bf Make clear what estimates you would consider as final if you could only choose on result to report.} 
Overall, I expect that 3 pages should suffice; 5 pages is a maximum.

You can work on this project in groups of 2 or 3 people. It is recommended that you split the work, {\bf and then give feedback on each other's report}. If you split the work, the project should not take too much work, but otherwise it may. In choosing a way to split the work, note that questions 1 and 5 don't require a lot of work (as they do not require analysis of the data), whereas question 4 may be perceived as a more difficult question.
One report per group must be returned no later than no later than May 15, together with structured software code (i.e. a script file in R or Python). All the files that you submit, must carry the family name of one of the members in the group. 

Make sure your report briefly mentions who did what. One third of the score will be on your own contribution and two thirds on the project as a whole (as indeed you are also responsible for the entirety of the project).

\end{document}
\begin{thebibliography}{111}
\bibitem[1]{Luque2018} Luque-Fernandez MA, Schomaker M, Rachet B, Schnitzer ME. Targeted maximum likelihood estimation for a binary treatment: A tutorial. Statistics in Medicine. 2018;37:2530-2546. 
\end{thebibliography}


